\documentclass[14pt,a4paper]{extarticle}
% =============================================================
% preamble.tex — оформление отчётов по СТП БГУИР 01-2024
% =============================================================

% Шрифты: под pdflatex — Times через tempora + T2A,
% под xelatex/tectonic — системный Times New Roman (запасной вариант STIX Two Text)
\usepackage{iftex}
\ifPDFTeX
  \usepackage[utf8]{inputenc}
  \usepackage[T2A]{fontenc}
  \usepackage{tempora}
\else
  \usepackage{fontspec}
  \IfFontExistsTF{Times New Roman}
    {\setmainfont{Times New Roman}}
    {\setmainfont{STIX Two Text}}
  \IfFontExistsTF{DejaVu Sans Mono}
    {\setmonofont{DejaVu Sans Mono}[Scale=0.82]}
    {\setmonofont{DejaVuSansMono}[Extension=.ttf, BoldFont=*-Bold,
                                  ItalicFont=*-Oblique, Scale=0.82]}
\fi
\usepackage[russian]{babel}

% Поля: слева 30 мм, справа 15 мм, сверху и снизу 20 мм
\usepackage{geometry}
\geometry{left=30mm, right=15mm, top=20mm, bottom=20mm}

\usepackage{setspace}
\singlespacing

\pretolerance=1000
\tolerance=2000
\emergencystretch=3em

\usepackage{indentfirst}
\setlength{\parindent}{1.25cm}

% Номер страницы — снизу справа
\usepackage{fancyhdr}
\pagestyle{fancy}
\fancyhf{}
\fancyfoot[R]{\thepage}
\renewcommand{\headrulewidth}{0pt}
\fancypagestyle{plain}{%
  \fancyhf{}%
  \fancyfoot[R]{\thepage}%
  \renewcommand{\headrulewidth}{0pt}%
}

\usepackage{amsmath,amssymb}
\usepackage{graphicx}
\usepackage{float}
\usepackage{booktabs}
\usepackage{tabularx}
\usepackage{multirow}
\usepackage{longtable}
\usepackage{array}

\usepackage{caption}
\captionsetup[figure]{labelsep=endash, justification=centering,
                      font={normalsize}, name=Рисунок}
\captionsetup[table]{labelsep=endash, justification=raggedright,
                     singlelinecheck=false, font={normalsize}, name=Таблица}

\usepackage{titlesec}
\titleformat{\section}
  {\newpage\normalfont\fontsize{14}{18}\selectfont\bfseries}
  {\thesection}{0.5em}{\MakeUppercase}
\titlespacing*{\section}{1.25cm}{0pt}{14pt}
\titleformat{\subsection}
  {\normalfont\fontsize{14}{18}\selectfont\bfseries}
  {\thesubsection}{0.5em}{}
\titlespacing*{\subsection}{1.25cm}{14pt}{14pt}
\titleformat{\subsubsection}
  {\normalfont\fontsize{14}{18}\selectfont\bfseries}
  {\thesubsubsection}{0.5em}{}
\titlespacing*{\subsubsection}{1.25cm}{14pt}{14pt}

\newcommand{\sectioncenter}[1]{%
  \newpage
  \phantomsection
  \addcontentsline{toc}{section}{#1}
  \begin{center}\textbf{\MakeUppercase{#1}}\end{center}
  \vspace{14pt}
}

\usepackage{tocloft}
\renewcommand{\cftsecfont}{\normalfont}
\renewcommand{\cftsecpagefont}{\normalfont}
\renewcommand{\cftsubsecfont}{\normalfont}
\renewcommand{\cftsubsecpagefont}{\normalfont}
\renewcommand{\cftsecleader}{\cftdotfill{\cftdotsep}}
\renewcommand{\cftdotsep}{1}
% babel переопределяет заголовки после преамбулы, поэтому правим через \addto
\addto\captionsrussian{\renewcommand{\contentsname}{\hfill\textbf{СОДЕРЖАНИЕ}\hfill}}
\renewcommand{\cftaftertoctitle}{\hfill}

\usepackage{enumitem}
\setlist[itemize]{noitemsep, topsep=0pt, leftmargin=1.25cm, label=--}
\setlist[enumerate]{noitemsep, topsep=0pt, leftmargin=1.25cm}

% Листинги кода и вывода программы
\usepackage{listings}
\usepackage{fancyvrb}
\lstset{
  basicstyle=\ttfamily\footnotesize,
  frame=single,
  breaklines=true,
  showstringspaces=false,
  inputencoding=utf8,
  extendedchars=true,
  columns=fullflexible,
  keepspaces=true,
  captionpos=t
}
\ifPDFTeX
\lstset{
  literate={а}{{\cyra}}1 {б}{{\cyrb}}1 {в}{{\cyrv}}1 {г}{{\cyrg}}1
           {д}{{\cyrd}}1 {е}{{\cyre}}1 {ж}{{\cyrzh}}1 {з}{{\cyrz}}1
           {и}{{\cyri}}1 {й}{{\cyrishrt}}1 {к}{{\cyrk}}1 {л}{{\cyrl}}1
           {м}{{\cyrm}}1 {н}{{\cyrn}}1 {о}{{\cyro}}1 {п}{{\cyrp}}1
           {р}{{\cyrr}}1 {с}{{\cyrs}}1 {т}{{\cyrt}}1 {у}{{\cyru}}1
           {ф}{{\cyrf}}1 {х}{{\cyrh}}1 {ц}{{\cyrc}}1 {ч}{{\cyrch}}1
           {ш}{{\cyrsh}}1 {щ}{{\cyrshch}}1 {ъ}{{\cyrhrdsn}}1
           {ы}{{\cyrery}}1 {ь}{{\cyrsftsn}}1 {э}{{\cyrerev}}1
           {ю}{{\cyryu}}1 {я}{{\cyrya}}1 {ё}{{\cyryo}}1
           {А}{{\CYRA}}1 {Б}{{\CYRB}}1 {В}{{\CYRV}}1 {Г}{{\CYRG}}1
           {Д}{{\CYRD}}1 {Е}{{\CYRE}}1 {Ж}{{\CYRZH}}1 {З}{{\CYRZ}}1
           {И}{{\CYRI}}1 {Й}{{\CYRISHRT}}1 {К}{{\CYRK}}1 {Л}{{\CYRL}}1
           {М}{{\CYRM}}1 {Н}{{\CYRN}}1 {О}{{\CYRO}}1 {П}{{\CYRP}}1
           {Р}{{\CYRR}}1 {С}{{\CYRS}}1 {Т}{{\CYRT}}1 {У}{{\CYRU}}1
           {Ф}{{\CYRF}}1 {Х}{{\CYRH}}1 {Ц}{{\CYRC}}1 {Ч}{{\CYRCH}}1
           {Ш}{{\CYRSH}}1 {Щ}{{\CYRSHCH}}1 {Ъ}{{\CYRHRDSN}}1
           {Ы}{{\CYRERY}}1 {Ь}{{\CYRSFTSN}}1 {Э}{{\CYREREV}}1
           {Ю}{{\CYRYU}}1 {Я}{{\CYRYA}}1 {Ё}{{\CYRYO}}1
           {—}{{---}}1 {–}{{--}}1 {…}{{\ldots}}1
}
\fi

% Подписи листингов — в том же стиле, что таблицы: «Листинг X.Y — Название»
\captionsetup[lstlisting]{labelsep=endash, justification=raggedright,
                          singlelinecheck=false, font={normalsize}}
\renewcommand{\lstlistingname}{Листинг}
\AtBeginDocument{\numberwithin{lstlisting}{section}}

\lstdefinestyle{python}{language=Python, keywordstyle=\bfseries}
\lstdefinestyle{shell}{language=bash, keywordstyle=\mdseries}
\lstdefinestyle{plain}{language={}, basicstyle=\ttfamily\scriptsize}

% Абзац после листинга должен сохранять отступ 1,25 см (СТП 2.1.1)
\makeatletter
\AtBeginDocument{\let\@endpetrue\@endpefalse}
\makeatother

% Подпись к блоку verbatim: \captionof вне плавающего окружения
% глобально обнуляет \parindent, поэтому отступ восстанавливается вручную
\newcommand{\listingcaption}[2]{%
  \captionof{lstlisting}{#1}\label{#2}%
  \setlength{\parindent}{1.25cm}%
}

\usepackage{hyperref}
\hypersetup{colorlinks=true, linkcolor=black, urlcolor=blue, citecolor=black}

\numberwithin{figure}{section}
\numberwithin{table}{section}
\numberwithin{equation}{section}

% Титульный лист отчёта по лабораторной работе
\newcommand{\makelabtitle}[2]{%
  \begin{titlepage}
  \thispagestyle{empty}
  \begin{center}

  МИНИСТЕРСТВО ОБРАЗОВАНИЯ РЕСПУБЛИКИ БЕЛАРУСЬ

  \vspace{0.5em}
  Учреждение образования

  \vspace{0.5em}
  \textbf{<<Белорусский государственный университет информатики и радиоэлектроники>>}

  \vspace{2em}
  Факультет компьютерных систем и сетей

  \vspace{0.5em}
  Кафедра программного обеспечения информационных технологий
  \end{center}

  \vspace{3em}
  \begin{center}
  \textbf{ОТЧЕТ}

  \vspace{0.5em}
  по лабораторной работе №#1

  \vspace{0.5em}
  по дисциплине

  \vspace{0.5em}
  <<Информационные и технические средства защиты авторских прав>>

  \vspace{0.5em}
  <<#2>>
  \end{center}

  \vspace{5em}
  \begin{flushright}
  \begin{minipage}{0.5\textwidth}
  \textbf{Выполнил:}\\
  Лебедевич Артем Владимирович\\
  магистрант группы 556301

  \vspace{0.6em}
  \textbf{Преподаватель:}\\
  Ярмолик Вячеслав Николаевич\\
  доктор технических наук\\
  профессор кафедры ПОИТ
  \end{minipage}
  \end{flushright}

  \vfill
  \begin{center}Минск, 2026~г.\end{center}
  \end{titlepage}
}

% Титульный лист реферата
\newcommand{\makerefattitle}[1]{%
  \begin{titlepage}
  \thispagestyle{empty}
  \begin{center}

  МИНИСТЕРСТВО ОБРАЗОВАНИЯ РЕСПУБЛИКИ БЕЛАРУСЬ

  \vspace{0.5em}
  Учреждение образования

  \vspace{0.5em}
  \textbf{<<Белорусский государственный университет информатики и радиоэлектроники>>}

  \vspace{2em}
  Факультет компьютерных систем и сетей

  \vspace{0.5em}
  Кафедра программного обеспечения информационных технологий
  \end{center}

  \vspace{3em}
  \begin{center}
  \textbf{РЕФЕРАТ}

  \vspace{0.5em}
  по дисциплине

  \vspace{0.5em}
  <<Информационные и технические средства защиты авторских прав>>

  \vspace{1.5em}
  на тему

  \vspace{0.5em}
  \textbf{<<#1>>}
  \end{center}

  \vspace{5em}
  \begin{flushright}
  \begin{minipage}{0.5\textwidth}
  \textbf{Выполнил:}\\
  Лебедевич Артем Владимирович\\
  магистрант группы 556301

  \vspace{0.6em}
  \textbf{Преподаватель:}\\
  Ярмолик Вячеслав Николаевич\\
  доктор технических наук\\
  профессор кафедры ПОИТ
  \end{minipage}
  \end{flushright}

  \vfill
  \begin{center}Минск, 2026~г.\end{center}
  \end{titlepage}
}


\begin{document}

\makelabtitle{1}{Шифрование методами перестановок и подстановок}

\tableofcontents

\section{Цель работы}

Изучить простейшие алгоритмы шифрования --- методы перестановок и методы подстановок.
Реализовать шифраторы и дешифраторы на основе трёх перестановочных методов
(<<железнодорожная изгородь>>, столбцовый метод, поворачивающаяся решётка), а также
подстановочных шифров по всем трём вариантам задания: шифр Цезаря, шифр Плейфейра и
шифр Виженера. Проверить корректность реализации на примерах из методического пособия.

\section{Теоретические сведения}

\subsection{Общая модель перестановочного шифрования}

При шифровании методами перестановок исходное сообщение $M$ делится на блоки по $d$
символов $M = m_1 \dots m_d,\ m_{d+1} \dots m_{2d},\ \dots$, после чего символы внутри
блока переставляются в соответствии с некоторым вектором, функцией или фигурой:
\begin{equation}
E_k(M) = m_{f(1)} \dots m_{f(d)},\ m_{f(d+1)} \dots m_{f(2d)},\ \dots
\end{equation}

Ключом служит сама перестановка $f$, дешифрование --- применение обратной
перестановки $f^{-1}$, поэтому шифр является симметричным. Принципиальное свойство
перестановки: символы не заменяются, а лишь меняются местами, следовательно частоты
букв шифротекста в точности совпадают с частотами букв исходного текста. Это делает
перестановочные шифры уязвимыми к криптоанализу, но одновременно полезными в качестве
слоя диффузии в составе современных блочных шифров.

Перед шифрованием текст нормализуется: отбрасываются пробелы и знаки препинания,
буквы приводятся к верхнему регистру, буква <<Ё>> отождествляется с <<Е>>.
Шифротекст принято выводить группами по пять символов.

\subsection{Метод <<железнодорожная изгородь>>}

Символы исходного текста записываются в виде, напоминающем по форме забор
(зигзагом по $K$ строкам), а символы шифротекста считываются из полученной записи
построчно. Высота изгороди $K$ является ключом. Номер строки для позиции $i$
определяется периодом $P = 2K - 2$:
\begin{equation}
r(i) =
\begin{cases}
p, & p < K,\\
P - p, & p \geqslant K,
\end{cases}
\qquad p = i \bmod P.
\end{equation}

Для $M = \text{<<Это лабораторная работа по КиОКИ>>}$ и $K = 4$ запись имеет вид,
приведённый в таблице~\ref{tab:fence}, а шифротекст получается чтением строк сверху вниз:
$C = \text{<<ЭОНОИ ТБРРА БТКОО ААОЯА АОКЛТ РПИ>>}$.

\begin{table}[H]
\caption{Запись текста зигзагом при $K = 4$}
\label{tab:fence}
\centering
\begin{tabular}{ll}
\toprule
Строка & Символы \\
\midrule
0 & Э \quad О \quad Н \quad О \quad И \\
1 & Т \; Б \; Р \; Р \; А \; Б \; Т \; К \; О \\
2 & О \; А \; А \; О \; Я \; А \; А \; О \; К \\
3 & Л \quad Т \quad Р \quad П \quad И \\
\bottomrule
\end{tabular}
\end{table}

Число различных ключей равно числу допустимых значений $K$, то есть шифр вскрывается
полным перебором высоты изгороди за время, линейное по длине текста.

\subsection{Столбцовый метод}

Исходное сообщение записывается в таблицу построчно, а затем считывается оттуда
столбцами согласно вектору, задающему порядок считывания. Вектор задаётся ключевым
словом: буквам ключа назначаются номера в соответствии с алфавитом, а если буква
встречается несколько раз, нумерация определяется порядком следования повторяющихся
букв в ключевом слове. Для ключа $K = \text{<<КРИПТОГРАФИЯ>>}$ получаем вектор
$(5, 8, 3, 7, 10, 6, 2, 9, 1, 11, 4, 12)$, а таблица шифрования приведена в
таблице~\ref{tab:column}.

\begin{table}[H]
\caption{Шифрование столбцовым методом с ключом <<КРИПТОГРАФИЯ>>}
\label{tab:column}
\centering
\begin{tabular}{cccccccccccc}
\toprule
5 & 8 & 3 & 7 & 10 & 6 & 2 & 9 & 1 & 11 & 4 & 12 \\
\midrule
Э & Т & О & Л & А & Б & О & Р & А & Т & О & Р \\
Н & А & Я & Р & А & Б & О & Т & А & П & О & К \\
И & О & К & И &   &   &   &   &   &    &   &   \\
\bottomrule
\end{tabular}
\end{table}

Считывая столбцы в порядке возрастания номеров, получаем
$C = \text{<<ААООО ЯКООЭ НИББЛ РИТАО РТААТ ПРК>>}$. Число ключей равно $d!$, где
$d$~--- длина ключевого слова, однако шифр вскрывается методом анаграммирования:
столбцы сопоставляются между собой по частотам биграмм русского языка.

Во время Первой мировой войны немецкие военные использовали \emph{двойной столбцовый
шифр}, при котором одно сообщение шифровалось дважды с одним и тем же или с двумя
разными ключами. Для второго ключа $K_2 = \text{<<ШИФРОВАНИЕ>>}$ с вектором
$(10, 4, 9, 8, 7, 2, 1, 6, 5, 3)$ получается
$C = \text{<<КИРЯР ПЭОАИ ТОАОТ КОЛТО БАОБА АНР>>}$. Двойная перестановка разрушает
столбцы первой таблицы и существенно усложняет анаграммирование, однако и она была
вскрыта: французской разведке требовалось несколько дней после смены ключа.

\subsection{Метод поворачивающейся решётки}

Исходный текст записывается на бумаге через отверстия решётки, которая по мере
заполнения поворачивается на 90$^\circ$. Решётка строится так: в матрице $n \times n$
($n$~--- чётное) ячейкам, занимающим при повороте на 90$^\circ$ одинаковое положение,
присваивается один и тот же номер, после чего вырезается по одному квадрату для
каждого номера. При повороте клетка $(r, c)$ переходит в клетку
\begin{equation}
(r, c) \longrightarrow (c,\ n - 1 - r),
\end{equation}
то есть все клетки матрицы разбиваются на $n^2/4$ орбит по четыре клетки. Выбор
ровно одной клетки из каждой орбиты гарантирует, что за четыре поворота решётка
покроет каждую клетку матрицы ровно один раз.

Разметка матрицы $4 \times 4$ и использованная решётка приведены в
таблице~\ref{tab:grille}. Пустые клетки в конце заполняются ничего не значащими
буквами.

\begin{table}[H]
\caption{Разметка орбит поворота и отверстия решётки $4 \times 4$}
\label{tab:grille}
\centering
\begin{tabular}{cccc@{\hskip 2cm}cccc}
\toprule
\multicolumn{4}{c@{\hskip 2cm}}{Разметка} & \multicolumn{4}{c}{Отверстия} \\
\midrule
1 & 2 & 3 & 1 & $\blacksquare$ & $\square$ & $\square$ & $\square$ \\
3 & 4 & 4 & 2 & $\square$ & $\square$ & $\square$ & $\blacksquare$ \\
2 & 4 & 4 & 3 & $\square$ & $\square$ & $\blacksquare$ & $\square$ \\
1 & 3 & 2 & 1 & $\square$ & $\blacksquare$ & $\square$ & $\square$ \\
\bottomrule
\end{tabular}
\end{table}

Для $M = \text{<<ШИФРОТЕКСТ>>}$ порядок заполнения матрицы по четырём поворотам
приведён в таблице~\ref{tab:grille-fill}; жирным выделены буквы-заполнители.

\begin{table}[H]
\caption{Заполнение матрицы по поворотам решётки}
\label{tab:grille-fill}
\centering
\begin{tabular}{lll}
\toprule
Поворот & Клетки $(r, c)$ & Записывается \\
\midrule
0$^\circ$ & (0,0) (1,3) (2,2) (3,1) & Ш И Ф Р \\
90$^\circ$ & (0,3) (1,0) (2,1) (3,2) & О Т Е К \\
180$^\circ$ & (0,2) (1,1) (2,0) (3,3) & С Т \textbf{А Б} \\
270$^\circ$ & (0,1) (1,2) (2,3) (3,0) & \textbf{В Г Д Е} \\
\bottomrule
\end{tabular}
\end{table}

Итоговая матрица считывается построчно, что даёт
$C = \text{<<ШВСОТ ТГИАЕ ФДЕРК Б>>}$. Число ключей для решётки $n \times n$
составляет $4^{n^2/4}$: для решётки $4 \times 4$ это всего 256 вариантов, а для
решётки $8 \times 8$ --- уже $4^{16} \approx 4{,}3 \cdot 10^{9}$.

\subsection{Шифр Цезаря и аффинное преобразование}

Подстановочные шифры, в отличие от перестановочных, заменяют символы, сохраняя их
позиции. Шифр Цезаря заменяет каждый символ алфавита исходного текста на сдвинутый
относительно него на $k$ позиций символ того же алфавита:
\begin{equation}
E_k(i) = (i + k) \bmod n, \qquad D_k(i) = (i + n - k) \bmod n.
\end{equation}

Обобщением является метод децимации $E_k(i) = (i \cdot k) \bmod n$ и объединяющее оба
метода аффинное преобразование с двумя ключами:
\begin{equation}
E_k(i) = (i \cdot k_1 + k_2) \bmod n, \qquad
D_k(j) = k_1^{-1} (j - k_2) \bmod n.
\label{eq:affine}
\end{equation}

Числа $k_1$ и $n$ должны быть взаимно простыми: в противном случае обратного элемента
$k_1^{-1}$ не существует и разные буквы исходного текста склеиваются в одну, то есть
преобразование перестаёт быть обратимым. В работе реализовано именно
преобразование~(\ref{eq:affine}); классический шифр Цезаря получается при $k_1 = 1$.
Мощность алфавита составляет $n = 26$ для английского языка и $n = 32$ для русского
(буква <<Ё>> отождествлена с <<Е>>).

\subsection{Шифр Плейфейра}

Биграммный шифр Плейфейра применялся в Великобритании во время Первой мировой войны и
является наиболее известным полиграммным шифром замены. Основой шифра служит шифрующая
таблица, которая представляет собой сеансовый ключ; для удобства запоминания она
заполняется ключевым словом, а затем оставшимися буквами алфавита. Для английского
языка используется таблица $5 \times 5$, в которой позиция для символа <<J>>
соответствует позиции символа <<I>> (таблица~\ref{tab:playfair}).

\begin{table}[H]
\caption{Шифрующая таблица для ключа <<CRYPTOGRAPHY>>}
\label{tab:playfair}
\centering
\begin{tabular}{|c|c|c|c|c|}
\hline
C & R & Y & P & T \\ \hline
O & G & A & H & B \\ \hline
D & E & F & I/J & K \\ \hline
L & M & N & Q & S \\ \hline
U & V & W & X & Z \\ \hline
\end{tabular}
\end{table}

Исходный текст разбивается на биграммы $m_i m_{i+1}$, которые преобразуются по
правилам:
\begin{enumerate}
\item буквы в одной строке --- берутся два символа справа от них (первый столбец
считается столбцом справа по отношению к последнему): $\text{YT} \to \text{PC}$;
\item буквы в одном столбце --- берутся символы ниже (первая строка считается строкой
ниже последней): $\text{VG} \to \text{RE}$;
\item буквы в разных строках и столбцах --- берутся два других угла прямоугольника,
причём $c_i$ находится в той же строке, что и $m_i$: $\text{ZF} \to \text{WK}$;
\item при $m_i = m_{i+1}$ между ними вставляется пустой символ <<X>>:
$\text{SS} \to \text{SXS} \to \text{QZ}$;
\item при нечётной длине текста пустой символ добавляется в конец.
\end{enumerate}

Для $M = \text{<<CIPHERTEXT>>}$ шифрование даёт $C = \text{<<PDHIMGRKZP>>}$.
Шифрование биграммами существенно повышает стойкость по сравнению с простой заменой,
поскольку выравнивает частотные характеристики, однако частотные свойства
распределения биграмм по-прежнему остаются ключом для злоумышленника.

\subsection{Шифр Виженера}

Шифр Виженера является многоалфавитным шифром подстановки и развивает идею Цезаря:
таблицы подстановки чередуются в зависимости от последовательности символов ключа.
Каждая строка таблицы Виженера представляет собой шифр Цезаря для конкретного
значения ключа, поэтому шифрование описывается выражением
\begin{equation}
E_k(i) = (m_i + k_{i \bmod d}) \bmod n,
\end{equation}
где $d$~--- длина ключа. Различают три варианта использования ключа:
\begin{enumerate}
\item \textbf{прямое использование} --- ключ циклически повторяется;
\item \textbf{прогрессивный ключ} --- при повторном применении символы ключа
циклически сдвигаются на одну позицию в упорядоченном алфавите: <<MODE>>,
<<NPEF>>, <<OQFG>> и так далее;
\item \textbf{самогенерирующийся ключ} --- в качестве последующих символов ключа
используется сам исходный текст.
\end{enumerate}

Результаты шифрования слова $M = \text{<<cryptography>>}$ ключом <<MODE>> по всем трём
схемам приведены в таблице~\ref{tab:vigenere}.

\begin{table}[H]
\caption{Шифр Виженера при различных схемах применения ключа}
\label{tab:vigenere}
\centering
\begin{tabularx}{\textwidth}{lXl}
\toprule
Схема & Последовательность ключа & Шифротекст \\
\midrule
Прямая & MODEMODEMODE & OFBTF CJVMD KC \\
Прогрессивная & MODENPEFOQFG & OFBTG DKWOF ME \\
Самогенерирующаяся & MODECRYPTOGR & OFBTV FEGTD NP \\
\bottomrule
\end{tabularx}
\end{table}

\section{Реализация}

\subsection{Структура программы}

Программа написана на языке Python~3 без внешних зависимостей и разделена на модули
по методам шифрования (листинг~\ref{lst:tree}).

\listingcaption{Структура проекта}{lst:tree}
\begin{Verbatim}[fontsize=\footnotesize, frame=single, samepage=true]
main.py              CLI: выбор метода, режима и ключа
configs/cfg.py       ключи и параметры по умолчанию
helpers/text.py      нормализация текста и вывод группами по пять символов
helpers/keyword.py   числовой ключ (порядок столбцов) из ключевого слова
helpers/alphabet.py  выбор рабочего алфавита: русский (32) или английский (26)
lab1/fence.py        железнодорожная изгородь
lab1/columnar.py     столбцовый метод и двойная перестановка
lab1/grille.py       поворачивающаяся решётка
lab1/caesar.py       шифр Цезаря и аффинное преобразование
lab1/playfair.py     биграммный шифр Плейфейра
lab1/vigenere.py     шифр Виженера (три схемы ключа)
lab1/runner.py       прогон всех примеров методички с проверкой
gui/methods.py       описание шифров для окна: параметры, вызовы, разбор
gui/app.py           виджеты: выбор шифра, поля параметров, решётка
\end{Verbatim}

Во всех шифрах дешифратор не дублирует логику шифратора, а переиспользует расчёт
позиций или сдвигов, выполняя обратное преобразование. Это гарантирует согласованность
пары <<шифрование--дешифрование>> и вдвое сокращает объём кода.

\subsection{Ключевые фрагменты кода}

Расчёт номера строки зигзага (листинг~\ref{lst:fence}) является общей основой
шифрования и дешифрования для метода изгороди: шифратор читает символы по строкам,
дешифратор раскладывает их обратно по тем же позициям.

\begin{lstlisting}[style=python, caption={Метод <<железнодорожная изгородь>>}, label={lst:fence}]
def row_pattern(length: int, key: int) -> List[int]:
    """Номер строки для каждой позиции текста: 0,1,..,K-1,K-2,..,1,0,1,..."""
    if key < 2:
        return [0] * length
    period = 2 * key - 2
    return [p if p < key else period - p for p in (i % period for i in range(length))]


def encrypt(text: str, key: int) -> str:
    rows = row_pattern(len(text), key)
    return "".join(ch for row in range(key)
                   for ch, r in zip(text, rows) if r == row)
\end{lstlisting}

Числовой ключ столбцового метода получается сортировкой позиций букв ключевого слова
по паре <<буква, позиция>>: вторичный ключ сортировки обеспечивает требуемую нумерацию
повторяющихся букв слева направо (листинг~\ref{lst:keyword}).

\begin{lstlisting}[style=python, caption={Порядок чтения столбцов из ключевого слова}, label={lst:keyword}]
def read_order(keyword: str) -> List[int]:
    letters = normalize(keyword)
    return sorted(range(len(letters)), key=lambda i: (letters[i], i))
\end{lstlisting}

При дешифровании столбцового метода необходимо учитывать, что последняя строка таблицы
может быть неполной: первые $|C| \bmod d$ столбцов длиннее остальных на один символ.
Без этого учёта весь текст сдвигается и восстановление становится невозможным
(листинг~\ref{lst:columnar}).

\begin{lstlisting}[style=python, caption={Восстановление длин столбцов при дешифровании}, label={lst:columnar}]
    full, tail = divmod(len(cipher), width)
    lengths = [full + (1 if c < tail else 0) for c in range(width)]
\end{lstlisting}

Для решётки ключ проверяется на корректность через разметку орбит: если два отверстия
попали в одну орбиту, при повороте они совпадут и часть матрицы останется незаполненной
(листинг~\ref{lst:grille}).

\begin{lstlisting}[style=python, caption={Разметка орбит поворота и порядок заполнения}, label={lst:grille}]
def rotate(cell: Cell, n: int) -> Cell:
    """Поворот клетки на 90 градусов по часовой стрелке."""
    row, col = cell
    return col, n - 1 - row


def fill_order(holes: List[Cell], n: int) -> List[int]:
    """Позиции матрицы в порядке заполнения: 4 поворота, внутри - построчно."""
    order: List[int] = []
    current = list(holes)
    for _ in range(4):
        order.extend(row * n + col for row, col in sorted(current))
        current = [rotate(cell, n) for cell in current]
    return order
\end{lstlisting}

Аффинное преобразование реализовано одной функцией: при дешифровании множитель
обращается по модулю мощности алфавита встроенной функцией \texttt{pow(mul, -1, n)}
(листинг~\ref{lst:caesar}).

\begin{lstlisting}[style=python, caption={Шифр Цезаря как частный случай аффинного преобразования}, label={lst:caesar}]
    if decode:
        inverse = pow(mul, -1, n)
        code = lambda i: (inverse * (i - add)) % n
    else:
        code = lambda i: (i * mul + add) % n
\end{lstlisting}

Реализация шифра Плейфейра не привязана к таблице $5 \times 5$: размеры вычисляются
из мощности алфавита, поэтому для русского языка автоматически строится таблица
$4 \times 8$. Три правила шифрования и дешифрования различаются только направлением
сдвига, что задаётся параметром \texttt{step}~$= \pm 1$ (листинг~\ref{lst:playfair}).

\begin{lstlisting}[style=python, caption={Правила шифра Плейфейра}, label={lst:playfair}]
        (r1, c1), (r2, c2) = position[pair[0]], position[pair[1]]
        if r1 == r2:                        # одна строка - сдвиг вбок
            c1, c2 = (c1 + step) % columns, (c2 + step) % columns
        elif c1 == c2:                      # один столбец - сдвиг вниз
            r1, r2 = (r1 + step) % rows, (r2 + step) % rows
        else:                               # прямоугольник - обмен столбцами
            c1, c2 = c2, c1
\end{lstlisting}

Три схемы применения ключа Виженера сведены к одной функции расчёта сдвига
(листинг~\ref{lst:vigenere}). Особенность самогенерирующегося ключа в том, что при
дешифровании ключ восстанавливается по ходу работы: очередной его символ берётся из
уже расшифрованной части текста, поэтому дешифратор передаёт в функцию накопленный
результат, а шифратор --- исходный текст.

\begin{lstlisting}[style=python, caption={Три схемы применения ключа Виженера}, label={lst:vigenere}]
def _key_shift(i: int, key: Sequence[int], mode: str,
               plain: Sequence[str], alphabet: str) -> int:
    """Сдвиг для позиции i; для autokey используется уже известный текст."""
    length = len(key)
    if mode == "autokey":
        return key[i] if i < length else alphabet.index(plain[i - length])
    step = i // length if mode == "progressive" else 0
    return (key[i % length] + step) % len(alphabet)
\end{lstlisting}

\subsection{Интерфейс командной строки}

Программа управляется из командной строки: первый аргумент задаёт метод, второй ---
режим (\texttt{enc}~--- шифрование, \texttt{dec}~--- дешифрование), далее указываются
ключ и текст. Текст может быть передан флагом \texttt{-t} или через стандартный ввод,
флаг \texttt{-{}-raw} отключает разбиение вывода на группы по пять символов
(листинг~\ref{lst:cli}).

\begin{lstlisting}[style=shell, caption={Примеры запуска}, label={lst:cli}]
python3 main.py demo

python3 main.py fence    enc -k 4 -t "Это лабораторная работа по КиОКИ"
python3 main.py column   enc -k КРИПТОГРАФИЯ -t "Это лабораторная работа по КиОКИ"
python3 main.py column   enc -k КРИПТОГРАФИЯ -d ШИФРОВАНИЕ -t "Это лабораторная..."
python3 main.py grille   enc -k 0,7,10,13 -t "ШИФРОТЕКСТ"

python3 main.py caesar   enc -k 3 -t "Это лабораторная работа по КиОКИ"
python3 main.py caesar   enc -k 5 -m 3 -t "Cryptography"
python3 main.py playfair enc -k CRYPTOGRAPHY -t "ciphertext"
python3 main.py vigenere enc -k MODE -s progressive -t "Cryptography"

echo "ШВСОТТГИАЕФДЕРКБ" | python3 main.py grille dec
\end{lstlisting}

Некорректные ключи отсекаются с понятным сообщением: перекрывающиеся при повороте
отверстия решётки, множитель аффинного шифра, не взаимно простой с мощностью алфавита,
символы ключа вне рабочего алфавита.


\subsection{Графическое приложение}

Помимо командной строки реализовано окно на библиотеке \texttt{tkinter} (входит в
стандартную поставку Python, дополнительных зависимостей не требуется), запускаемое
командой \texttt{python3 main.py gui}. Приложение построено так, чтобы параметры
шифрования были открыты для изменения, а шифры выбирались по одному, а не применялись
все сразу.

Окно состоит из четырёх частей:
\begin{enumerate}
\item выбор шифра из списка и общие настройки вывода (размер группы символов);
\item панель параметров, которая перестраивается под выбранный шифр;
\item два поля текста --- открытый текст и шифротекст --- с кнопками прямого и
обратного преобразования;
\item панель <<Как это работает>>, показывающая промежуточное представление:
зигзаг изгороди, таблицу столбцов с номерами, разметку орбит и порядок заполнения
решётки, шифрующую таблицу Плейфейра или последовательность ключа Виженера под текстом.
\end{enumerate}

Чтобы интерфейс не приходилось править при добавлении нового шифра, описание шифров
вынесено в отдельный модуль: каждый шифр задаётся записью \texttt{Method} со списком
параметров \texttt{Param} и функциями шифрования, дешифрования и разбора
(листинг~\ref{lst:gui}). Виджет для параметра выбирается по его типу
(\texttt{int}~--- поле с шагом, \texttt{text}~--- строка ввода, \texttt{choice}~---
выпадающий список, \texttt{grille}~--- интерактивная решётка), поэтому добавление
шифра сводится к одной записи в словаре \texttt{METHODS}.

\begin{lstlisting}[style=python, caption={Описание шифра для интерфейса}, label={lst:gui}]
    "Шифр Виженера": Method(
        title="Подстановка: каждая позиция шифруется своим сдвигом из ключа",
        params=[
            Param("keyword", "Ключевое слово", "text", "MODE"),
            Param("scheme", "Схема ключа", "choice", "direct", choices=vigenere.MODES),
        ],
        encrypt=lambda text, v: vigenere.encrypt(text, v["keyword"], v["scheme"]),
        decrypt=lambda text, v: vigenere.decrypt(text, v["keyword"], v["scheme"]),
        explain=_explain_vigenere,
    ),
\end{lstlisting}

Отдельно решена задача ввода ключа поворачивающейся решётки: вместо списка номеров
клетки выбираются мышью, а клик по клетке автоматически снимает выбор с ранее выбранной
клетки той же орбиты. Благодаря этому ключ остаётся корректным по построению --- в нём
всегда ровно одно отверстие на орбиту, и невозможно задать решётку, которая оставит
часть матрицы незаполненной.

\section{Результаты работы программы}

Команда \texttt{python3 main.py demo} выполняет все реализованные шифры на примерах из
методического пособия и сравнивает полученные шифротексты с эталонными. Сводные
результаты приведены в таблице~\ref{tab:results}, полный вывод программы ---
в листинге~\ref{lst:demo}.

\begin{table}[H]
\caption{Результаты шифрования}
\label{tab:results}
\centering
\begin{tabularx}{\textwidth}{lXl}
\toprule
Метод & Ключ & Совпадение \\
\midrule
Железнодорожная изгородь & $K = 4$ & да \\
Столбцовый метод & КРИПТОГРАФИЯ & да \\
Двойная перестановка & КРИПТОГРАФИЯ, ШИФРОВАНИЕ & да \\
Поворачивающаяся решётка & $4 \times 4$, отверстия 0, 7, 10, 13 & да \\
Шифр Цезаря & $k = 3$ & эталона нет \\
Аффинное преобразование & $k_1 = 3$, $k_2 = 5$ & эталона нет \\
Шифр Плейфейра & CRYPTOGRAPHY & да \\
Виженер, прямой ключ & MODE & да \\
Виженер, прогрессивный ключ & MODE & да \\
Виженер, самогенерирующийся ключ & MODE & да \\
\bottomrule
\end{tabularx}
\end{table}

Все восемь шифротекстов, для которых методическое пособие приводит эталонное значение,
совпали с ним посимвольно. Дешифрование во всех случаях восстанавливает исходный текст
с точностью до служебных символов --- букв-заполнителей решётки и разделителей
Плейфейра, которые получатель отбрасывает по смыслу сообщения.

Дополнительно выполнена проверка обратимости на случайных текстах: 300 текстов длиной
до 60 символов для каждого перестановочного метода при всех допустимых значениях ключа
и 400 текстов для подстановочных методов при всех сдвигах, всех допустимых множителях
аффинного шифра и всех трёх схемах ключа Виженера. Расхождений не зафиксировано.

\listingcaption{Вывод программы}{lst:demo}
\VerbatimInput[fontsize=\scriptsize, frame=single]{listings/demo_output.txt}

\section{Сравнение методов и стойкость}

Сводные характеристики реализованных шифров приведены в таблице~\ref{tab:compare},
где $N$~--- длина текста, $d$~--- длина ключевого слова, $n$~--- мощность алфавита.

\begin{table}[H]
\caption{Сравнение реализованных шифров}
\label{tab:compare}
\centering
\begin{tabularx}{\textwidth}{lllX}
\toprule
Шифр & Ключей & Сложность & Способ вскрытия \\
\midrule
Изгородь & $N$ & $O(N)$ & полный перебор высоты \\
Столбцовый & $d!$ & $O(N)$ & анаграммирование по биграммам \\
Двойной столбцовый & $(d!)^2$ & $O(N)$ & анаграммирование, требует больше текста \\
Решётка & $4^{n^2/4}$ & $O(N)$ & перебор при малом $n$, статистика текста \\
Цезарь & $n$ & $O(N)$ & полный перебор сдвига \\
Аффинный & $n \cdot \varphi(n)$ & $O(N)$ & частотный анализ, перебор \\
Плейфейр & $25!$ & $O(N)$ & частотный анализ биграмм \\
Виженер & $n^d$ & $O(N)$ & тест Казиски, индекс совпадений \\
\bottomrule
\end{tabularx}
\end{table}

Перестановочные шифры полностью сохраняют частоты одиночных букв, поэтому уже сам факт
совпадения частотного профиля шифротекста с профилем языка указывает криптоаналитику
на применённый класс шифра. Подстановочные шифры, наоборот, сохраняют позиции символов:
одноалфавитная замена вскрывается частотным анализом, биграммный Плейфейр --- анализом
частот биграмм, а многоалфавитный Виженер сначала сводится к $d$ независимым шифрам
Цезаря определением периода ключа.

\sectioncenter{Заключение}

В ходе лабораторной работы изучены простейшие алгоритмы шифрования и реализованы на
языке Python шифраторы и дешифраторы трёх перестановочных методов
(<<железнодорожная изгородь>>, столбцовый метод с одинарным и двойным ключом,
поворачивающаяся решётка) и трёх подстановочных шифров, соответствующих всем вариантам
задания (Цезарь с обобщением до аффинного преобразования, Плейфейр, Виженер с тремя
схемами применения ключа). Корректность подтверждена совпадением с восемью эталонными
примерами методического пособия и проверками обратимости на случайных данных.

Практический вывод работы состоит в том, что ни перестановка, ни подстановка по
отдельности не обеспечивают криптостойкости: первая сохраняет частоты символов, вторая
--- их позиции, и обе вскрываются статистическими методами при доступном объёме
шифротекста в несколько сотен символов. Значение этих преобразований в другом: в
современных блочных шифрах (DES, AES) они применяются совместно и многораундово, где
подстановка обеспечивает конфузию --- нелинейную связь ключа и шифротекста, а
перестановка обеспечивает диффузию --- распространение влияния одного бита открытого
текста на весь блок.

\sectioncenter{Список использованных источников}

\begin{enumerate}
\item Ярмолик, С.~В. Криптография и охрана коммерческой информации: метод. пособие по
выполнению лаб. работ / С.~В.~Ярмолик, В.~Н.~Ярмолик. --- Минск: БГУИР, 2010. --- 32~с.
\item Шнайер, Б. Прикладная криптография. Протоколы, алгоритмы, исходные тексты на
языке Си / Б.~Шнайер. --- М.: Триумф, 2002. --- 816~с.
\item Столлингс, В. Криптография и защита сетей: принципы и практика /
В.~Столлингс. --- М.: Вильямс, 2001. --- 672~с.
\end{enumerate}

\end{document}
